\section{Transport, observable variance, and trace mixing}
\label{sec:transport}

Fix an annealing parameter $j$ and suppose that $p\geq 1$.
Write $f=f_j$, $\widetilde f=\widetilde f_j$, and $C=C(j)$.
We use the paired ground set, transversal class $\mathcal A$, defect
classes $\mathcal A_{il}$, and their conditional masses from
Section~\ref{sec:signatures}. In particular, $f$ is positive on
\emph{all} $p$-subsets of the $2p$ elements, and has the omission
property. The second lift $\widetilde f$ has the addition property.
They agree on transversals and satisfy the comparison
\eqref{eq:lift-comparison}, with $q=W^2$, on the state space
\[
  \mathcal X=\mathcal A\,\cup\!
       \bigcup_{\substack{1\leq i,l\leq p\\i\ne l}}\mathcal A_{il}.
\]
We give signed-flow, observable-variance, and trace arguments for these
two lifts, keeping track of the comparison factors and of the domains
on which they are needed. The flow-energy and trace viewpoints are
classical; see Aldous and Fill~\cite[Sections~2.7.1 and~3.7.1]{AF2002}.
The transport inequality first bounds transversal variance and defect
means; together these control the particular trajectory averages used
by the algorithm. The transversal variance bound also gives mixing
of the transversal trace, whose return-time analysis bounds its cost.

We adapt the paired-transversal transport and observable/trace analysis
for counting common bases of two matroids
in~\cite[Sections~4--6]{OpenAICommonBases2026}.
Here the two required signature properties belong to separate lifts.
Their comparison contributes the factors $q^2$ and $q^3$ in the
transversal-variance and defect-mean bounds below. We give the complete
argument for these weights.

\subsection{The enlarged Metropolis chain}

For each ordered pair $i\ne l$, choose a positive multiplier $w_{il}$,
and put
\[
 \lambda(S)=
 \begin{cases}
   f(S),&S\in\mathcal A,\\
   w_{il}f(S),&S\in\mathcal A_{il}.
 \end{cases}
 \qquad
 \Lambda=\sum_{S\in\mathcal X}\lambda(S),
 \qquad
 \pi(S)=\frac{\lambda(S)}{\Lambda}.
\]
Call the multipliers \emph{good} if
\begin{equation}
 \frac{C}{4c_{il}}\leq w_{il}\leq\frac{4C}{c_{il}}
 \qquad(i\ne l).
 \label{eq:good-multipliers}
\end{equation}
For good multipliers,
\begin{equation}
 \Lambda\leq \bigl(1+4p(p-1)\bigr)C\leq 5p^2C,
 \qquad
 \pi(\mathcal A)\geq\frac{1}{5p^2},
 \qquad
 \pi(\mathcal A_{il})\geq\frac{1}{20p^2}.
 \label{eq:type-mass-bounds}
\end{equation}
Indeed, a defect type has total multiplied mass between $C/4$ and
$4C$, and the transversal mass is $C$. When $p=1$ there are no defect
types; all assertions concerning them are omitted.

The chain holds with probability $1/2$. Otherwise it independently
chooses one uniform element of $S$ and one uniform element of its
complement and proposes to exchange them. A proposal outside
$\mathcal X$ is rejected. A proposal $S'$ in $\mathcal X$ is accepted
with probability
$\min\{1,\lambda(S')/\lambda(S)\}$.
Denote its transition operator by $P$.
Every defect state can reach a transversal by moving one element of
its full pair to its empty pair. Transversals are connected by flips
within individual pairs. Positivity of $\lambda$ therefore makes the
chain irreducible. Its holding probability is at least $1/2$.
For distinct exchange neighbors,
\[
 \pi(S)P(S,S')=
 \frac{\min\{\lambda(S),\lambda(S')\}}{2p^2\Lambda},
\]
so $P$ is reversible with stationary law $\pi$.
For a real function $g$ on $\mathcal X$, define
\begin{align}
 \mathfrak D(g)
   &=\sum_{\{S,S'\}}
       \min\{\lambda(S),\lambda(S')\}
          \bigl(g(S)-g(S')\bigr)^2,
 \notag\\
 \mathcal E(g)
   &=\langle g,(I-P)g\rangle_\pi
     =\frac{\mathfrak D(g)}{2p^2\Lambda}.
 \label{eq:energy}
\end{align}
The sum is over unordered exchange edges in $\mathcal X$.
This convention will be used throughout the section.

\subsection{A transport inequality on one-hole states}

The next lemma uses the omission property of $f$ on its full domain.
Some configurations contributing to its auxiliary masses need not
belong to $\mathcal X$. Its applications will identify the one-hole states and
their conductances with actual states and conductances of the
Metropolis chain.

\begin{lemma}[One-hole transport]
\label{lem:transport}
Fix disjoint sets of included and excluded elements. Let $F$ be the
set of fixed inclusions, and partition the remaining elements into two
singleton slots $P_0,P_1$ and $s$ ordinary slots of two elements each.
Assume $|F|+s+1=p$.
Let $\mathcal X_{\mathrm{sl}}$ consist of the sets that respect the
fixed assignments and select one element from every slot except one
hole slot, from which they select nothing.
Give a state with hole slot $i$ weight $W_i f(S)$, where every $W_i$ is
positive. For a real function $g$ on $\mathcal X_{\mathrm{sl}}$, let
$D_{\mathrm{sl}}(g)$ be the unordered-edge energy on its single-exchange
graph, with conductance equal to the minimum of the two endpoint weights.

Let $s_0,s_1$ be the unmultiplied $f$-masses of the states with hole
$P_0,P_1$, and let $\bar g_0,\bar g_1$ be their respective $f$-weighted
means. For each ordinary slot $t$, let $v_t$ be the unmultiplied mass
of sets that omit both singleton slots, select both elements of $t$,
and select one element from every other ordinary slot, while
respecting the fixed assignments. Then
\begin{equation}
 (\bar g_0-\bar g_1)^2
 \leq
 \left(
   \frac{1}{s_0W_{P_0}}+\frac{1}{s_1W_{P_1}}
   +\frac{2}{s_0s_1}\sum_t\frac{v_t}{W_t}
 \right)D_{\mathrm{sl}}(g).
 \label{eq:slot-transport}
\end{equation}
The notation $W_i$ here denotes slot multipliers, independently of the
capacity threshold $W$.
\end{lemma}

\begin{proof}
We seek a signed edge current whose divergence is $f(S)/s_0$ at
states with hole $P_0$, $-f(S)/s_1$ at states with hole $P_1$, and zero
at all ordinary-hole states. Orient each edge once, write $I_{SS'}$
for its signed current, and write $c_{SS'}$ for its conductance.
Such a current satisfies, by summation by parts and Cauchy--Schwarz,
\begin{equation}
 \begin{split}
  (\bar g_0-\bar g_1)^2
   &=\left(\sum_{\{S,S'\}}
                 I_{SS'}\bigl(g(S)-g(S')\bigr)\right)^2\\
   &\leq
      \left(\sum_{\{S,S'\}}\frac{I_{SS'}^2}{c_{SS'}}\right)
      D_{\mathrm{sl}}(g).
 \end{split}
 \label{eq:transport-flow-criterion}
\end{equation}
It therefore suffices to bound its flow energy by the coefficient in
\eqref{eq:slot-transport}.

The local routing step explains the energy quantity we will use.
A complete display $R$ contains the two singleton elements and one
chosen element from every ordinary slot. Its one-hole sets
$S_i=F\cup(R\setminus\{i\})$, $i\in R$, form a clique of exchange
edges. For an element $i$, write $W_i$ for its slot's multiplier.
Put $m_i=f(S_i)$. For any coefficients $h_i$ with
$\sum_i m_i h_i=0$, regard $m_i h_i$ as the prescribed signed
divergence at $S_i$. Choose a hub $h_*$ maximizing $W_i m_i$, and
send current $m_i h_i$ from each nonhub $S_i$ to the hub. Balance gives
the prescribed divergence at the hub as well. Each used edge has
conductance $W_i m_i$ at its nonhub endpoint, so this routing has
energy at most
\[
 \sum_{i\ne h_*}\frac{(m_i h_i)^2}{W_i m_i}
 \leq\sum_{i\in R}\frac{m_i h_i^2}{W_i}.
\]
We now assign these balanced demands recursively. Their sum over
complete displays will have the required source and sink divergences;
the recursion will also bound the sum of their routing energies.

Fix an order of the ordinary slots. Initially display the two
singleton elements. Expose ordinary slots in that order, branching
into two children according to which element of the next slot is
added to the display.
At a node write $R$ for its displayed elements, excluding $F$.
For $i\in R$, let $m_i$ be the sum of $f$ over sets that contain
$F\cup(R\setminus\{i\})$, omit the nondisplayed mates in already
processed slots, and have single occupancy in every unprocessed
ordinary slot. These sets all have size $p$, and $m_i>0$.
At a complete display this definition reduces to $m_i=f(S_i)$.
Maintain coefficients $h_i$, interpreted as demand per unit
unmultiplied weight, satisfying
\begin{equation}
  \sum_{i\in R}m_i h_i=0.
  \label{eq:transport-balance}
\end{equation}
At the root their values are $1/s_0$ and $-1/s_1$ at the elements of
$P_0$ and $P_1$, respectively. Assign the node the potential
\[
  \Psi(R,h)=\sum_{i\in R}\frac{m_i h_i^2}{W_i}.
\]
At a leaf this potential bounds the routing energy just computed. Each $W_i$
is fixed by its slot throughout the recursion. These multipliers
enter the energy bound and the eventual hub choice, but not the
recursive coefficients or prescribed divergences.

Suppose the next slot $t$ has elements $a,b$.
For an old hole $i\in R$, let $m_i^a,m_i^b$ be its masses in the two
children. Then $m_i^a+m_i^b=m_i$.
The mass of the newly introduced hole, at $a$ in the first child and
at $b$ in the second, is the same positive number $u$: in either case
the entire slot $t$ is empty.
Put $X=\sum_{i\in R}h_i m_i^a$.
Retain every old coefficient unchanged, and give the new hole
coefficient $-X/u$ in the $a$ child and $X/u$ in the $b$ child.
Equation~\eqref{eq:transport-balance} holds in both children because
$\sum_i h_i m_i^b=-X$.
The sum of their potentials exceeds the parent potential by exactly
\begin{equation}
  \frac{2X^2}{uW_t}.
  \label{eq:transport-increment}
\end{equation}

We bound these increments using the omission signature. At any node,
and for any unprocessed ordinary slot $t$, define a symmetric matrix
$V_t(R)$ on the displayed elements, with zero diagonal. For distinct
$i,k\in R$, its entry is the sum of $f$ over configurations with holes
at $i,k$ in the display, double occupancy in slot $t$, and single
occupancy in every other unprocessed slot. Fixed inclusions,
exclusions, and previously excluded mates remain in force.
All these configurations have size $p$.

When $t=\{a,b\}$ is the next slot, $V_t(R)$ is the old block of
\[
 N=\begin{pmatrix}
     V_t(R)&(m_i^b)_{i\in R}&(m_i^a)_{i\in R}\\
     (m_i^b)_{i\in R}^{\mathsf T}&0&u\\
     (m_i^a)_{i\in R}^{\mathsf T}&u&0
   \end{pmatrix}.
\]
This form has at most one positive eigenvalue. To see the precise
application, fix the context and exclude the nondisplayed mates of
processed slots. Sum all unprocessed ordinary slots other than $t$
at occupancy one. On the remaining ground set $R\cup\{a,b\}$ the
resulting degree is $|R|$, so $N$ is its two-omission matrix.
Lemmas~\ref{lem:binary-signatures} and
\ref{lem:binary-contraction} give the asserted signature.
In particular, this argument uses the full-domain weight $f$ even
when an entry of $V_t(R)$ describes multiple empty and full original
pairs. No comparison with $\widetilde f$ is made here.

Extend $h$ by zero at $a,b$ and put $d_{ab}=e_a-e_b$.
The vector $e_a+e_b$ has square $2u>0$ under $N$.
Both $h$ and $d_{ab}$ are orthogonal to it in this form, since
\eqref{eq:transport-balance} holds and the two new diagonal entries
are zero. A symmetric form with at most one positive eigenvalue is
nonpositive on the form-orthogonal complement of a positive-square
vector: another positive-square vector in that complement would span
a positive definite plane with the first vector. On the complement,
Cauchy--Schwarz for the negative of the form applies. Here
\[
 h^{\mathsf T}N d_{ab}=-2X,
 \qquad d_{ab}^{\mathsf T}N d_{ab}=-2u,
 \qquad h^{\mathsf T}Nh=h^{\mathsf T}V_t(R)h.
\]
Consequently
\begin{equation}
   \frac{2X^2}{u}\leq -h^{\mathsf T}V_t(R)h.
   \label{eq:transport-schur}
\end{equation}

Fix now a slot $t$ to be processed at a later level.
When a preceding slot is split, the sum of the two child quantities
$h^{\mathsf T}V_t(R)h$ equals the parent quantity.
The old entries split additively. A cross entry involving the new
hole represents a configuration in which the split slot is empty,
so it is identical in the two children. Its contributions cancel
because the new coefficients are opposite and all old coefficients
are retained. There is no new diagonal contribution.
Iterating this identity shows that, immediately before processing
$t$, the sum of these quadratic quantities over the whole level
is their value at the root. The root matrix has off-diagonal entry
$v_t$ and zero diagonal, so that value is
$-2v_t/(s_0s_1)$.
By \eqref{eq:transport-increment} and
\eqref{eq:transport-schur}, the total potential increase at the level
of slot $t$ is at most $2v_t/(s_0s_1W_t)$.
The root potential is
$1/(s_0W_{P_0})+1/(s_1W_{P_1})$.
Thus the sum of all leaf potentials is bounded by the parenthetical
coefficient in \eqref{eq:slot-transport}.

Route the prescribed demands at every leaf by the clique construction
above. Its energy is at most $\Psi(R,h)$.
Different leaves cannot use the same edge: the union of two distinct
endpoints of a leaf edge is exactly $F\cup R$, and so determines the
entire display and its unique sequence of choices. Hence the assembled
flow has energy at most the sum of the leaf potentials.

It remains to identify the divergence after summing the leaf flows.
A state with a singleton hole has one leaf representation; its
divergence is $f(S)/s_0$ or $-f(S)/s_1$.
A state $S$ with an ordinary hole at slot $t$ has exactly two
representations, obtained by adjoining either element of $t$ to
$S\setminus F$. Every other displayed ordinary choice is fixed by
$S$. The two histories agree before processing $t$, so the same
$X,u$ create opposite coefficients there. Those coefficients remain
unchanged at every later branch, regardless of the later masses and
new coefficients. At the leaves both masses are exactly $f(S)$,
and their divergences cancel. This reasoning also applies if the
state is a hub in either leaf, since its net divergence is still
$m_i h_i$.
The assembled flow therefore has the required divergence, and
\eqref{eq:transport-flow-criterion} together with the leaf-potential
bound proves \eqref{eq:slot-transport}. If $s=0$, the root is already
a complete display and the same clique argument applies directly.
\end{proof}

\subsection{Variance on transversals}

Assume henceforth that the multipliers are good, and write
\[
  \mu(S)=\frac{f(S)}{C}\quad(S\in\mathcal A),
  \qquad C_{\mathrm p}=p(1+8pq^2).
\]
The variance under $\mu$ of a function on $\mathcal X$ refers to its
restriction to $\mathcal A$.

\begin{proposition}
\label{prop:transversal-variance}
Suppose the multipliers satisfy~\eqref{eq:good-multipliers}.
For every real function $g$ on $\mathcal X$,
\begin{equation}
    C\Var_\mu(g)\leq C_{\mathrm p}\mathfrak D(g).
    \label{eq:transversal-variance}
\end{equation}
\end{proposition}

\begin{proof}
We decompose the transversal variance along a tree of partial
assignments, choosing the next pair as a function of the current
assignment. Let $\sigma$ fix some pairs to specified single choices.
For distinct unassigned labels define
\[
 a_{il}=
  \frac{\widetilde c_{il}^{\sigma}}{z_\sigma}
  \frac{\widetilde c_{li}}{C}.
\]
Lemma~\ref{lem:defect-inequalities} gives
\begin{equation}
 a_{il}a_{li}\leq\frac1{16},
 \qquad
 a_{il}a_{lk}\leq a_{ik}
 \quad\text{for distinct }i,l,k.
 \label{eq:sink-products}
\end{equation}
For the second inequality, use the conditional defect inequality
on the first factors and the unconditioned inequality with indices
$k,l,i$ on the second factors.
Form a directed graph with arc $i\to l$ when $a_{il}>1$.
It has no directed two-cycle. A shortest longer directed cycle could
be shortened using the second inequality in
\eqref{eq:sink-products}, also a contradiction.
Thus this graph has a sink $k$, for which $a_{kt}\leq1$ at every other
unassigned label $t$. If only one label remains, choose it as $k$.

Branch on the two choices in pair $k$.
Apply Lemma~\ref{lem:transport} with the two elements of $k$ as
singleton slots, both of multiplier one, and with the other
unassigned pairs as ordinary slots. Keep $\sigma$ fixed.
A singleton-hole state is a transversal respecting $\sigma$; an
ordinary-hole state at $t$ has type $tk$ and multiplier $W_t=w_{tk}$.
The auxiliary doubled-slot mass is exactly $v_t=c_{kt}^\sigma$.
Thus every one-hole state in this application belongs to
$\mathcal X$ with precisely its weight $\lambda$.
The good-multiplier bound, followed by two uses of
\eqref{eq:lift-comparison}, gives
\begin{align}
 \frac{v_t}{W_t}
  =\frac{c_{kt}^\sigma}{w_{tk}}
  &\leq\frac{4c_{kt}^\sigma c_{tk}}{C}
   \leq\frac{4q^2\widetilde c_{kt}^\sigma
                        \widetilde c_{tk}}{C}
 \notag\\
  &=4q^2z_\sigma a_{kt}\leq4q^2z_\sigma.
 \label{eq:sink-comparison}
\end{align}
These comparisons involve only actual defect types.

Let $s_0,s_1$ and $\bar g_0,\bar g_1$ be the masses and means of the
two transversal children, so $s_0+s_1=z_\sigma$.
Their unnormalized between-means variance contribution is
\[
   \frac{s_0s_1}{z_\sigma}(\bar g_0-\bar g_1)^2.
\]
Let $\mathfrak D_\sigma(g)$ be the sum in
\eqref{eq:energy} restricted to edges whose two endpoints respect
$\sigma$. The slot energy in this application is a subgraph energy,
so it is at most $\mathfrak D_\sigma(g)$.
Multiplying \eqref{eq:slot-transport} by $s_0s_1/z_\sigma$ and using
\eqref{eq:sink-comparison} yields
\begin{equation}
 \frac{s_0s_1}{z_\sigma}(\bar g_0-\bar g_1)^2
 \leq (1+8pq^2)\mathfrak D_\sigma(g).
 \label{eq:conditional-variance-charge}
\end{equation}
The first two terms in the transport coefficient contribute exactly
$(s_0+s_1)/z_\sigma=1$; the ordinary slots contribute at most
$8pq^2$.

Iterate this branching until every pair is fixed. The conditional
variance identity decomposes $C\Var_\mu(g)$ as the sum of the
between-means contributions over the internal nodes.
At any depth, two distinct nodes disagree on the choice of the pair
fixed at their first divergence. Their restricted state sets are
disjoint, even if subsequent pairs were chosen in different orders.
Consequently their internal edge sets are disjoint and
\[
  \sum_{\sigma\text{ at a fixed depth}}\mathfrak D_\sigma(g)
  \leq\mathfrak D(g).
\]
An edge crossing a fixed-pair choice is internal to neither of the
two resulting restrictions; it does not create a duplicate charge.
There are $p$ possible branching depths. Summing
\eqref{eq:conditional-variance-charge} proves
\eqref{eq:transversal-variance}.
The conditioning tree and sink choices are used only in this proof.
\end{proof}

\subsection{Defect means and observable variance}

Set
\[
  C_{\mathrm d}=8+32pq^3.
\]
For $i\ne l$, let
$\bar g_{il}=c_{il}^{-1}\sum_{S\in\mathcal A_{il}}f(S)g(S)$.

\begin{lemma}
\label{lem:defect-mean}
Suppose the multipliers satisfy~\eqref{eq:good-multipliers}, and let
$g$ be a real function on $\mathcal X$.
For each defect type $il$ there is a conditional transversal mean
$\bar g'$ such that
\begin{equation}
 (\bar g_{il}-\bar g')^2
       \leq \frac{C_{\mathrm d}}{C}\mathfrak D(g),
 \qquad
 (\bar g'-\mathbb E_\mu g)^2\leq4\Var_\mu(g).
 \label{eq:defect-mean-bounds}
\end{equation}
\end{lemma}

\begin{proof}
Of the four possible joint choices in pairs $i,l$, choose one whose
transversal mass is at least $C/4$.
Denote the selected elements by $a_i,a_l$ and their mates by
$b_i',b_l'$.
In Lemma~\ref{lem:transport}, fix $a_l$ included and $b_i'$ excluded,
and take $P_0=\{a_i\}$, $P_1=\{b_l'\}$. All other pairs are ordinary
slots. A hole at $P_0$ gives the entire defect type $il$, whereas a
hole at $P_1$ gives the selected conditional transversal class.
Thus
\[
 s_0=c_{il},\quad W_{P_0}=w_{il},\qquad
 s_1\geq C/4,\quad W_{P_1}=1.
\]
An ordinary hole at $t$ has type $tl$, restricted to the choice $a_i$
in pair $i$, and has multiplier $W_t=w_{tl}$.
A doubled ordinary slot with both singleton slots empty is a subset
of type $it$, so $v_t\leq c_{it}$.
All one-hole states and their exchange edges are therefore genuine
states and edges of the chain, with the same endpoint weights.
Their energy is at most $\mathfrak D(g)$.

The lift comparison and Lemma~\ref{lem:defect-inequalities} give,
for each ordinary slot $t$,
\begin{equation}
 c_{it}c_{tl}
 \leq q^2\widetilde c_{it}\widetilde c_{tl}
 \leq q^2C\widetilde c_{il}
 \leq q^3C c_{il}.
 \label{eq:defect-comparison}
\end{equation}
The first two terms in the coefficient of
\eqref{eq:slot-transport} are at most $4/C+4/C$.
For its last term, \eqref{eq:good-multipliers} and
\eqref{eq:defect-comparison} give
\[
 \frac{2}{s_0s_1}\sum_t\frac{v_t}{W_t}
 \leq\frac{8}{c_{il}s_1C}\sum_t c_{it}c_{tl}
 \leq\frac{8pq^3}{s_1}
 \leq\frac{32pq^3}{C}.
\]
This proves the first inequality of
\eqref{eq:defect-mean-bounds} with $\bar g'$ the selected conditional
mean. The conditioning event has $\mu$-probability at least $1/4$;
conditional Cauchy--Schwarz gives
\[
 (\bar g'-\mathbb E_\mu g)^2
 \leq \mathbb E_\mu\bigl[(g-\mathbb E_\mu g)^2
                        \mid\text{selected choices}\bigr]
 \leq4\Var_\mu(g),
\]
which proves the second inequality.
\end{proof}

Define $\mathcal Q$ on real functions on $\mathcal X$ by retaining
all values on $\mathcal A$ and replacing the values on each defect
type by their mean within that type. Since its multiplier is
constant, the conditional $\pi$-mean on $\mathcal A_{il}$ is
$\bar g_{il}$. Hence $\mathcal Q$ is the orthogonal projection in
$L^2(\pi)$ associated with the partition into individual transversal
states and entire defect types. Its range consists of functions that
are arbitrary on $\mathcal A$ and constant within each defect type.
This includes type indicators and every function supported on
$\mathcal A$, the two kinds of observables used in
Section~\ref{sec:algorithm}. It preserves constants and means.

Chen, Vigoda, and Yang use a restricted variance bound for hole-pattern
frequencies in permanent approximation~\cite[Theorem~5 and
Lemma~12]{CVY2026}. Their partition pools the perfect matchings;
here each transversal is retained separately, while each defect type
is averaged. The inequality below controls this projection without
requiring a bound on all fluctuations within a defect type.
Put
\begin{equation}
  K_{\mathrm{obs}}=10p^4(9C_{\mathrm p}+2C_{\mathrm d}).
  \label{eq:observable-constant}
\end{equation}

\begin{proposition}[Observable variance]
\label{prop:observable-variance}
Suppose the multipliers satisfy~\eqref{eq:good-multipliers}.
For every real $g$ on $\mathcal X$,
\begin{equation}
 \Var_\pi(\mathcal Qg)
 \leq(9C_{\mathrm p}+2C_{\mathrm d})\frac{\mathfrak D(g)}{C}
 \leq K_{\mathrm{obs}}\mathcal E(g).
 \label{eq:observable-variance}
\end{equation}
\end{proposition}

\begin{proof}
Let $m=\mathbb E_\mu g$.
Lemma~\ref{lem:defect-mean} and $(a+b)^2\leq2a^2+2b^2$ give
\[
  (\bar g_{il}-m)^2
  \leq\frac{2C_{\mathrm d}}{C}\mathfrak D(g)
       +8\Var_\mu(g).
\]
Variance is no larger than squared deviation about any fixed
constant. Therefore
\begin{align*}
 \Var_\pi(\mathcal Qg)
 &\leq\mathbb E_\pi[(\mathcal Qg-m)^2]\\
 &=\pi(\mathcal A)\Var_\mu(g)
       +\sum_{i\ne l}\pi(\mathcal A_{il})(\bar g_{il}-m)^2\\
 &\leq9\Var_\mu(g)+\frac{2C_{\mathrm d}}{C}\mathfrak D(g)\\
 &\leq(9C_{\mathrm p}+2C_{\mathrm d})\frac{\mathfrak D(g)}{C},
\end{align*}
using Proposition~\ref{prop:transversal-variance} in the last step.
Finally, \eqref{eq:energy} and \eqref{eq:type-mass-bounds} give
$\mathfrak D(g)/C\leq10p^4\mathcal E(g)$.
When $p=1$ the defect sum is empty and $\mathcal Q=I$; the same
inequalities apply.
\end{proof}

\subsection{Trajectory averages of the controlled observables}

The covariance-sum viewpoint is classical; see
Aldous and Fill~\cite[Section~2.3]{AF2002}. Here we derive the needed
finite-time bound from the restricted observable inequality.

\begin{lemma}[Time averages from a dominated initial law]
\label{lem:time-average}
Suppose the multipliers satisfy~\eqref{eq:good-multipliers}.
Let $G$ be a real function on $\mathcal X$ with $\mathcal QG=G$.
Suppose $S_0$ has law $\nu$ satisfying $\nu(S)\leq U\pi(S)$ for
every state. For $N\geq1$ consecutive observations of the chain,
\begin{equation}
 \mathbb E_\nu\left[
    \left(\frac1N\sum_{t=0}^{N-1}G(S_t)-\mathbb E_\pi G\right)^2
                  \right]
 \leq\frac{2UK_{\mathrm{obs}}}{N}\Var_\pi(G).
 \label{eq:time-average-bound}
\end{equation}
\end{lemma}

\begin{proof}
Write $G_c=G-\mathbb E_\pi G$ and $v=\Var_\pi(G)$.
Because $\mathcal Q$ is an orthogonal projection fixing constants,
for every $g$,
\begin{align*}
 |\langle G_c,g\rangle_\pi|^2
 &=|\langle G_c,\mathcal Qg-\mathbb E_\pi g\rangle_\pi|^2\\
 &\leq v\Var_\pi(\mathcal Qg)
 \leq K_{\mathrm{obs}}v\mathcal E(g).
\end{align*}
On the mean-zero subspace, $I-P$ is positive definite and invertible:
reversibility makes $P$ self-adjoint, and finite irreducibility makes
its eigenvalue one simple. Take
$g=(I-P)^{-1}G_c$ on this subspace and put
\[
 a=\langle G_c,(I-P)^{-1}G_c\rangle_\pi=\mathcal E(g)\geq0.
\]
The previous inequality gives $a^2\leq K_{\mathrm{obs}}va$, and hence
$a\leq K_{\mathrm{obs}}v$, also when $a=0$.

Laziness puts the spectrum of $P$ in $[0,1]$: the operator $2P-I$ is
itself a reversible stochastic contraction. Thus the stationary
covariances
$c_t=\langle G_c,P^tG_c\rangle_\pi$ are nonnegative, and finite
spectral decomposition on the mean-zero subspace gives
\[
  \sum_{t=0}^{\infty}c_t
  =\langle G_c,(I-P)^{-1}G_c\rangle_\pi=a.
\]
For a stationary trajectory the mean squared error of the average is
\[
 \frac{Nv+2\sum_{t=1}^{N-1}(N-t)c_t}{N^2}
 \leq\frac{2a}{N}
 \leq\frac{2K_{\mathrm{obs}}v}{N}.
\]
Finally, the law of an entire trajectory started from $\nu$ is
bounded by $U$ times the corresponding stationary trajectory law:
only the initial factor changes in the product of transition
probabilities. Applying this domination to the nonnegative squared
error proves \eqref{eq:time-average-bound}.
\end{proof}

\subsection{The transversal trace and its return cost}

Starting from $x\in\mathcal A$, let
\[
  T_{\mathcal A}^+=\inf\{t\geq1:S_t\in\mathcal A\},
  \qquad
  P_{\mathrm{tr}}(x,y)
     =\mathbb P_x(S_{T_{\mathcal A}^+}=y).
\]
The trace makes one transition for each strictly positive return to
$\mathcal A$; consecutive visits and original holding steps count
as returns.
The mean return identity used below is the finite-state form of
Kac's formula~\cite[Corollary~2.24]{AF2002}.

\begin{lemma}[Trace mixing and return cost]
\label{lem:trace}
Suppose the multipliers satisfy~\eqref{eq:good-multipliers}.
The trace kernel $P_{\mathrm{tr}}$ is irreducible, reversible with
stationary law $\mu$, and has holding probability at least $1/2$.
Its inverse spectral gap is at most
$2p^2C_{\mathrm p}\leq K_{\mathrm{obs}}$.
In particular, with
$\mu_{\min}=\min_{x\in\mathcal A}\mu(x)>0$, for every integer $s\geq0$,
\begin{equation}
 \left|\frac{P_{\mathrm{tr}}^s(x,y)}{\mu(y)}-1\right|
 \leq\mu_{\min}^{-1}\exp(-s/K_{\mathrm{obs}})
 \qquad(x,y\in\mathcal A).
 \label{eq:trace-relative-mixing}
\end{equation}
If the initial trace law satisfies $\nu\leq U\mu$ pointwise, then
\begin{equation}
  \mathbb E_\nu T_{\mathcal A}^+
       \leq\frac{U}{\pi(\mathcal A)}.
  \label{eq:trace-return-cost}
\end{equation}
Furthermore, at every deterministic trace-step index the law remains
bounded by $U\mu$. The expected total number of underlying
transitions for $s$ trace steps is at most $sU/\pi(\mathcal A)$.
\end{lemma}

\begin{proof}
Finite irreducibility ensures finite expected hitting and return
times to $\mathcal A$. Explicitly, from each state choose a finite
positive-probability path to $\mathcal A$. Finiteness of the state
space gives a common bound on path length and a positive lower
bound on these path probabilities. The Markov property then gives
a geometric bound in blocks on the hitting-time tail. Applying this
after the first transition also bounds the strictly positive return
time. These qualitative bounds justify the harmonic extensions and
occupation sums below; no polynomial estimate is assumed at this
point.

For a function $h$ on $\mathcal A$, let $\widehat h$ be its harmonic
extension
\[
 \widehat h(x)=\mathbb E_x h(S_{T_{\mathcal A}}),
 \qquad T_{\mathcal A}=\inf\{t\geq0:S_t\in\mathcal A\}.
\]
Then $\widehat h=h$ on $\mathcal A$, $(I-P)\widehat h=0$ off
$\mathcal A$, and $(I-P)\widehat h=(I-P_{\mathrm{tr}})h$ on
$\mathcal A$ by conditioning on the first transition. For harmonic
extensions $\widehat h,\widehat k$ of $h,k$, it follows that
\begin{equation}
 \langle \widehat h,(I-P)\widehat k\rangle_\pi
   =\frac{C}{\Lambda}
        \langle h,(I-P_{\mathrm{tr}})k\rangle_\mu.
 \label{eq:trace-bilinear-identity}
\end{equation}
Symmetry of the left side proves reversibility of the stochastic
kernel $P_{\mathrm{tr}}$ with respect to $\mu$, and thus stationarity.
Tracing positive-probability paths of $P$ proves irreducibility.
Also $P_{\mathrm{tr}}(x,x)\geq P(x,x)\geq1/2$.

Let
$\mathcal E_{\mathrm{tr}}(h)
 =\langle h,(I-P_{\mathrm{tr}})h\rangle_\mu$.
Taking $k=h$ in \eqref{eq:trace-bilinear-identity} gives the exact
normalization
\[
  \mathfrak D(\widehat h)
    =2p^2\Lambda\mathcal E(\widehat h)
    =2p^2C\mathcal E_{\mathrm{tr}}(h).
\]
Proposition~\ref{prop:transversal-variance} therefore implies
\[
  \Var_\mu(h)\leq2p^2C_{\mathrm p}\mathcal E_{\mathrm{tr}}(h).
\]
This is the stated inverse-gap bound. The inequality
$2p^2C_{\mathrm p}\leq K_{\mathrm{obs}}$ follows immediately from
\eqref{eq:observable-constant} and $p\geq1$.

For the pointwise estimate put
$e_x(z)=\mathbf 1_{\{x\}}(z)/\mu(x)-1$.
Reversibility and spectral contraction on mean-zero functions give
\begin{align*}
 \left|\frac{P_{\mathrm{tr}}^s(x,y)}{\mu(y)}-1\right|
  &=|\langle e_x,P_{\mathrm{tr}}^s e_y\rangle_\mu|\\
  &\leq\exp(-s/K_{\mathrm{obs}})
        \sqrt{\bigl(\mu(x)^{-1}-1\bigr)
                    \bigl(\mu(y)^{-1}-1\bigr)}\\
  &\leq\mu_{\min}^{-1}\exp(-s/K_{\mathrm{obs}}),
\end{align*}
which proves \eqref{eq:trace-relative-mixing}.

To compute the mean return cost, start the underlying chain with law
$\mu$ on $\mathcal A$ and define the finite occupation measure
\[
  m(z)=\mathbb E_\mu
       \sum_{t=0}^{T_{\mathcal A}^+-1}\mathbf 1_{\{S_t=z\}}.
\]
Shifting the sum by one transition removes the initial law and adds
the return law. Both are $\mu$, since the trace is stationary.
Thus $mP=m$. Uniqueness of the stationary probability measure for
the finite irreducible chain implies $m=c\pi$ for some positive
constant $c$. Before the return, the only visit to $\mathcal A$ is at
time zero, so $m$ restricted to $\mathcal A$ is exactly $\mu$.
Consequently $c\pi(\mathcal A)=1$ and
\[
  \mathbb E_\mu T_{\mathcal A}^+
   =\sum_zm(z)=\frac{1}{\pi(\mathcal A)}.
\]
Pointwise domination of the starting law proves
\eqref{eq:trace-return-cost}.
Since $\mu$ is stationary for the trace,
$\nu P_{\mathrm{tr}}^r\leq U\mu$ for every deterministic integer
$r\geq0$. Applying \eqref{eq:trace-return-cost} to each of these
laws and using the Markov property at return times bounds the
expected sum of $s$ excursion lengths by $sU/\pi(\mathcal A)$.
This last bound is unconditional at the deterministic trace indices;
it does not assert domination after conditioning on an arbitrary
past event or on meeting a running-time cap.

When $p=1$, every state is a transversal, $P_{\mathrm{tr}}=P$, and
every return costs exactly one underlying transition. The stated
bounds include this case.
\end{proof}

The flows, auxiliary matrices, and adaptive conditioning trees in
this section certify the estimates but are not constructed by the
algorithm. Section~\ref{sec:algorithm} uses the directly simulated
Metropolis chain, the controlled observables, and capped trace
restarts, together with the bounded weight evaluator from
Proposition~\ref{prop:weight-oracle}.
